\section{Introduction}
\label{sec:intro}

Muon \citep{jordan2024muon} replaced the per-coordinate step of Adam-type
optimizers with a matrix-aware one: take the momentum of a weight matrix,
orthogonalize it with a few Newton--Schulz iterations, and step along the
result. Its preconditioned descendants go further and whiten the momentum
against the second moment of the layer's \emph{inputs} before orthogonalizing,
which is the construction we take as our baseline and call input-side MuonT
throughout. For a linear layer $y=Wx$, this uses the statistics of $x$.

The gradient of that layer, however, is an outer product of two vectors:
$G=\delta x^{\Tr}$, where $\delta$ is the cotangent arriving at the layer's
output. Input-side whitening conditions the update against the second moment
of $x$ and leaves the second moment of $\delta$ unused, even though $\delta$ is
already materialized in the backward pass and is no harder to sample than $x$
is. Kronecker-factored curvature approximations use both sides
\citep{martens2015kfac,gupta2018shampoo}, which suggests the output side
carries information the input side does not. This paper asks the narrow
empirical question that follows: if you add the output-cotangent second moment
to a Muon-type optimizer, what do you get, and what does it cost?

We answer it on a pinned recipe --- a 265M-parameter Llama-style transformer
trained on 2.68B tokens of a DCLM/StarCoder/Proof-Pile mixture, eight A100s per
run --- by defining a family of output-side normalized updates and running every
member of it to the full budget. The study comprises \nDistinctRuns{} distinct full-budget
trajectories across arm geometries, matched null controls, three seeds, two
batch sizes, and systems variants, plus 28 short-schedule learning-rate
screening trials. All of them train the same tokens in the same order within a seed,
which makes per-update paired differencing exact and is what lets us resolve
effects of a few thousandths of a nat at all.

\paragraph{What we find.}
Four results structure the paper, centered on \emph{width-thresholded output
normalization} (\mlpdiag{}, dense at width 512 and diagonal at width 2048) and
its pure-diagonal limit (\alldiag{}).

\begin{enumerate}
\item \textbf{The output side is a complement, not a substitute.}
  Added on top of input whitening, \mlpdiag{} improves held-out
  loss by \mlpdiagselmag{} nats and does so in every one of five training
  windows and on all three seeds we ran, while \alldiag{} retains two-thirds of
  that gain with no output eigendecompositions. Used \emph{instead} of input whitening,
  dense output factors make the optimizer \leftallsel{} nats
  \emph{worse} than the input-only baseline --- worse than either side alone
  (Section~\ref{sec:complement}).

\item \textbf{A matched null control absorbs a fifth to a third of the
  apparent gain.} Adding output-side machinery changes the numerical path, the
  communication pattern, and the host scheduling. Controls that execute all of
  that work and apply identity factors improve on the baseline by
  \nullsharedense\% of the dense arm's selection-loss gain and \nullsharemlpcfour\%
  of \mlpdiag{}'s C4 gain, leaving $-0.0059$ selection and $-0.0040$ C4 nats of
  net geometry benefit (Section~\ref{sec:null}).

\item \textbf{Hadamard decorrelation explains why width-thresholding works, and
  rules out step-size inflation.} Auditing all 352 checkpoint moments across 32
  layers shows that pre-SwiGLU width-2048 cotangents (\texttt{gate},
  \texttt{up}) have off-diagonal correlations at the finite-sample floor
  (median RMS \mechOutGateCorrRmsMed{} and \mechOutUpCorrRmsMed{}), so exact
  diagonals capture \mechOutGateExpVarLPct\% and \mechOutUpExpVarLPct\% of
  $\|L - I\|_F^2$ while taming $\mechOutGateDiagMaxMean\times$ coordinate-scale
  anisotropy; residual-stream width-512 cotangents (\texttt{o}, \texttt{down})
  are $97\%$ rotational, explaining why \mlpdiag{} beats both \dense{} and
  \alldiag{}. Simultaneously, instantaneous step norms match \muont{} to
  \mechMlpdiagGraftNorm{} ($\mechMlpdiagGraftNormPct\%$) while rotating updates
  by $\mechMlpdiagGraftDeg^\circ$, ruling out implicit step-size inflation
  (Section~\ref{sec:mechanism}).

\item \textbf{Lagged execution closes the latency gap to
  $+\sysFastMlpVsMuonLo\%$--$\sysFastMlpVsMuonHi\%$ over vanilla Muon.}
  Although the initial runner paid $+20.3\%$ over Muon ($+18.0\%$ over \muont{})
  for \texttt{inner\_outer\_mlpdiag\_lag}, stage profiling localized that gap to
  a per-step device-wide barrier, synchronous fallback stalls, 896 custom-op
  backward wrappers, and SIMT-FP32 matrix products under \texttt{allow\_tf32=False}.
  Removing those bottlenecks in \texttt{fast\_lag} cuts calibrated eight-A100
  step overhead over Muon to $+\sysFastMlpVsMuonLo\%$--$\sysFastMlpVsMuonHi\%$
  for \texttt{mlpdiag\_lag} ($+\sysFastMlpVsMuontLo\%$--$\sysFastMlpVsMuontHi\%$
  over \muont{}, saving \sysMemSavedGib\,GiB) and
  $+\sysFastAllVsMuonLo\%$--$\sysFastAllVsMuonHi\%$ even for fully dense
  \texttt{all\_lag}, crossing the \sysTokenSavingPct\% token-saving break-even
  line (Section~\ref{sec:systems}).
\end{enumerate}

\paragraph{What remains to confirm: the full-budget learning-rate horizon.}
One empirical gate remains before claiming tuned optimizer superiority. On a
short 320-update (6\% budget) screening schedule, \muont{} gained \lrHeadroom{}
nats when its learning rate moved from the recipe's $0.0226$ to $0.064$.
Carrying those short-schedule optima to the full 5{,}120-update budget,
however, \emph{degraded \muont{} itself by \waveTwelveMuontSelDiff{} nats}
relative to $0.0226$, and noisy short-horizon tie-breaking assigned $0.0905$ to
two diagonal candidates while assigning $0.064$ to a mathematically equivalent
fused arm, producing \transferWorst{}-nat regressions past the full-budget
stability cliff (Section~\ref{sec:tuning}). Thus $0.0226$, not $0.064$, lies
near the full-budget optimum, and Section~\ref{sec:program} specifies the
compact full-budget sweep around $0.0226$ on \texttt{fast\_lag}
\texttt{inner\_outer\_mlpdiag\_lag} and \alldiag{} that closes the comparison.

\paragraph{Scope.}
Everything here is one model size, one token budget, one hardware
configuration, and --- for the arm comparison --- a single data order per seed.
We make no claim about scaling: the clean statement of our results is about an
optimizer modification measured carefully at one point, not about a method.
Section~\ref{sec:limitations} is explicit about which of our numbers are
exploratory, which are replicated, and which are deterministic repeats of the
same trajectory rather than independent observations.
